\documentclass[10pt,norsk,a4paper,usenames,dvipsnames]{article}
\usepackage[utf8]{inputenc}
\usepackage[T1]{fontenc}
\usepackage[norsk]{babel}
% - PDF-relatert
\usepackage{hyperref,pdfpages,hypcap,download}
\hypersetup{colorlinks=true,allcolors=.}
\newcommand\fhref[2]{%
	\href{#1}{#2}\footnote{\url{#1}}%
}
% Andre pakker
\usepackage[cm]{fullpage}
\usepackage{parskip,multicol,textcomp,amssymb,graphicx,color,enumitem,cleveref}
% - Korreksjon av fotnoter i seksjoner/overskrifter
\usepackage[stable]{footmisc}
% - Skrifttype
\usepackage[bitstream-charter]{mathdesign}
% - Kommentarer
\usepackage{comment}
\usepackage{epstopdf}
\usepackage[dvipsnames]{xcolor}
\usepackage{xfrac}
% Redefines the first level of enumerated lists.
\renewcommand{\theenumi}{\alph{enumi}}
\renewcommand{\labelenumi}{\theenumi)}
\usepackage{setspace}


\epstopdfsetup{outdir=./}

\download[vedtekter.pdf]{protokoll/cyb-gf-protokoll-2021-vår.pdf}


\title{\huge Ordinær generalforsamling}
\author{\LARGE Cybernetisk Selskab}
\date{29. november 2022}

\begin{document}


\maketitle



\begin{center}


\includegraphics[width=0.6\textwidth,height=0.6\textheight,keepaspectratio=true]{cyblogoa3.pdf}

\end{center}


\newpage


\tableofcontents

\section{Valg av møteleder}
Chrsitina Eide velges ved akklamasjon.

\section{Valg av referent}
Iver Oknes velges ved akklamasjon.


\section{Valg av protokollunderskrivere}
Eric Reber og Martin Heggem velges ved akklamasjon.


\section{Valg av tellekorps}
Mons Eirik Bjerkan \& Mick Williams velges ved akklamasjon.

\section{Godkjenning av innkalling}
Innkalling godkjennes ved akklamasjon.


\section{Godkjenning av dagsorden}

\section{Semesterberetninger}
    \subsection{Semesterberetning ved leder}

        \begin{multicols}{2}

        Kjære alle sammen, takk for at dere har kunnet komme på generalforsamling idag. Dette semesteret har det vært mye nytt for mange i foreningen. Første vanlige fadderuke, Dagen@IFI og Pub til Pub siden 2019.

        \end{multicols}

        Semesterberetningen godkjennes ved akklamasjon.


    \subsection{Semesterberetning ved kjellermogul}

        \begin{multicols}{2}


            Dette semesteret har vært det første “normale” semesteret på flere år. Vi har fått til mye kult, og
            mye er nytt.

            Under fadderuka i år kunne vi endelig åpne uten antallsbegrensning. Siden ingen hadde
            erfaring med å arrangere fadderuke fra tidligere, ble ting gjort noe i siste liten, men vi syns
            resultatet ble bra. Vi hadde noe, men relativt lite problemer med å fylle opp skift, og det var både
            faddere og frivillige i CYB som holdt baren åpen hele uka. I forbindelse med fadderuka ble det
            også kjøpt inn nye samband som har fungert bra hele semesteret.

            Dagen@IFI ble også en stor suksess, mye grunnet tett arbeid mellom Dagen-styret og CYB.
            Generelt har vi sett dette semesteret at tett samarbeid med andre foreninger både gjør
            arrangementene bedre, men også bedrer CYBs posisjon i foreningslivet på IFI.

            I kaféen har vi i løpet av semesteret klart å fylle alle skiftene, som er ganske utrolig. Vi har fått
            nyheter, endret kaffekort, og mer annet. Espressomaskinen har fortsatt fungert (bank i bordet).

            Vi har også jobbet mye med økonomi. Store deler av dette arbeidet har vært gjort av
            økonomigruppa, primært økonomiansvarlig Hanne, og kasserer Trym, men jeg vil også rette en
            stor takk til Henrik Steen for dette.
            I kjellerstyret har vi sett på marginer og priser, og satt bedre rammer som sørger for at vi skal
            tjene nok penger for å drifte foreningen.
            Vi har også endret rutinene noe for utlån, som gjør det lettere å ta betalt, og gir litt mindre arbeid
            til utlånsansvarlig.

            Det nye kassesystemet har blitt godt tatt imot av alle både i og utenfor CYB, og har vært i bruk
            hele semesteret. Totalt siden innføringen, har vi i det nye kassesystemet omsatt for over 1
            million kroner.

            Vi mistet, etter forrige semester, våre DJ-sjefer. Det tok litt tid, men vi fikk til slutt etterfylt vervet,
            og jeg ser en lys fremtid for DJ-gruppa i CYB. Vi har fått innvilget penger av Kulturstyret til nytt
            DJ utstyr, som jeg tror kommer til å gjøre oss mer attraktive for både nye og erfarne DJer.

            I dette semesteret har vi hatt torsdagspuber og fredagspuber, og holdt åpent stort sett hver av
            dem, med noen få unntak. Vi har prøvd å legge samarbeidsarrangementer på barkvelder i
            stedet for som utlån. Dette har vi opplevd som veldig positivt for CYB, de andre foreningene og
            studentmassen på IFI. Vi har nesten fylt opp alle skift, med mange skjenkemestere og
            funksjonærer, mange av dem nye dette semesteret.

            I tillegg har vi dette semesteret, som forrige, hatt en flink kjellernestleder som sørger for at vi
            alltid har snacks og godt på møter.
            Jeg vil til slutt takke hele kjellerstyret for sitt gode arbeid, og et godt gjennomført semester, og
            jeg vil takke alle kaffemestere, skjenkemestere, arrmestere og funksjonærer som har stått på,
            og som gjør Escape til det flotte stedet det er for alle IFI-studenter.

        \end{multicols}

        Semesterberetningen godkjennes ved akklamasjon.


\section{Økonomiansvarlig orienterer om økonomi}
            % Feild pressentasjon presenteres, før riktig så represneteres.
            Økonomien er ryddet opp i, og det er sendt in mva. raporter tilbake til 2017.
            Økonomiansvarlig presenterer budsjett.
            Spørsmål fra dag:
                Det står i vedtektene at budsjett skal "behandles". Det konkluderes at dette betyr voteres over.
            Det voteres over budsjett:
                Budsjettet vedtas.


\section{Kontingentfastsettelse}
    Hovedstyret foreslår å holde medlemskontingenten på kr.~50,-.

    Fastsettelse godkjennes ved akklamasjon.


\section{Valg}
\begin{multicols}
    \subsection{Hovedstyret}

        Leder presenterer hva det vil si å ha verv i Hovedstyret.

        Det stilles spørsmål fra salen.

        Alle vervene som skal velges presenteres.

        Man velges inn i hovedstyret for to semestre av gangen.

        Det stilles spørsmål fra salen.

        % \subsubsection{Leder}


        % \subsubsection{Nestleder}


        \subsubsection{Kjellermogul}
        Det finnes en kandidat til Kjellermogul:
        Hanne Hugem:
        Går andre året prosa. Har vært intern i litt over et år. Vært ksjenkemester et esemseter, og jobbet i kafeen i tre semestere. Vært økonomiansavarlig nå, opg bestemt seg for å bli Kjellermogul nå.

        Spørsmål fra salen.

        \textbf{Hanne Alexandersdatter Hugem velges ved akklamasjon.}

        \subsubsection{Økonomiansvarlig}
        Det finnes en kandidat:
        Trym S. Auren
        Går tredjeåret robotikk, og er påutveksling i England dette året. Ønsker å bli Økonomiansvarlig får går vekk fra konto og over til prosjekt.

        \textbf{Trym Auren velges ved akklamasjon.}

        \subsubsection{Arrangementssjef}
        Det finnes en kandidat:
        Eric Reber
        Går tredjeåret bachelor på Prosa. Ikke vært med sålenge, men har likt seg godt her, så hvorfor ikke stille? Har tidligere arrangsjert mye forskjellig for andre foreninger.

        \textbf{Eric Emanuel Reber velges ved akklamasjon.}

        \subsubsection{Internansvarlig}
        Det finnes en kandidat.

        Kristian
        Elsker pingviner, og ønsker dermed å være internansvarlig.

        \textbf{Kristian Tereshchenko Stillingen velges ved akklamsjon.}

        \subsubsection{Promoteringssjef}
        Det finnes to kandidater til promoteringssjef.

        Fredrik
        Er nåværende kjellernestleder. Har jobbet med hvordan Cyb skal se ut utad fra foreningen. Ønsker dermed å bli promoteringsansvarlig for bedre kreativ kontroll på utad oppfatning.

        Ka Thas
        Er nestleder i Mikro, og har tidligere gått KDA. Har prøvd å være med i promoteringsgruppa tidligere, men fant en gruppe uten mye å gjøre. Ønsker å endre på dette.

        \textbf{Ka Thas velges ved alminnelig flertall.}

        \subsubsection{Rekrutteringsansvarlig}
        Det finnes en kandidater til Rekrutteringsansvarlig.

        Fredrik
        Synes han er veldig egnet for dette. Det synes også stedsfortreder.

        \textbf{Fredrik Sarai Velges ved akklamasjon.}


        \subsubsection{X-sjef}
        Det finens en kandidat til x-sjef.

        Kamil
        Har vært x-sjef 1 år nå. Går nest siste semester påmaster nå, så vil isåfall ta andre semester av vervet ute i jobb.

        \textbf{Kamil Krzy\żanowski velges ved akklamasjon.}


\end{multicols}
\begin{multicols}
    \subsection{Kjellerstyret} %TODO Oppdater listen med verv
        Alle verv som er til valg i kjellerstyret gjelder for ett semester av gangen. Med unntak av Økonomiansvarlig som blir valgt inn for to semester.

        Alle vervene som skal velges presenteres.

        \subsubsection{Kjellernestleder}
            % \textit{Vervet presenteres.}
            Det finnes en kandidat til vervet.

            Amalie Ingela Aaser

            Annser seg selv som ne venn av Hanne, og tenker at et sammarbeid mellom dem hadde fungert bra. Ønsker å få mere ting til å skje i kejlleren, så tenker å ved gode delegeringsegenskaper delegere mye. Ønsker også å være en motkraft til Hannes skepsis mot Jäger-tapp.

            \textbf{Amalie Ingela Aaser velges ved akklamasjon}



        \subsubsection{Barsjef}
            % \textit{Vervet presenteres.}
            Det finnes en kandidat til vervet.

            Mattis
            Sitter i styret i Mikro. Har Mye erfaring forran baren og litt bak. Mye kontakt med bryggerier i Østfold, og er deleier i Nøisom.

            \textbf{Mattis Stene-Johansen velges ved akklamasjon.}


        \subsubsection{Kafésjef}
            % \textit{Vervet presenteres.}
            Det finens en kandidat til vervet.

            Martin
            Har vært tilknyttet IFI lenge nok til å ha hatt INF1000 påpapir. Kan ikke såmye om kaffe, men er villig til å lære. Går første år prosa på master.

            \textbf{Martin Mihle Nygaard velges ved akklamasjon.}

        \subsubsection{Kasserer}
            % \textit{Vervet presenteres.}
            En kandidat til vervet

            Sima
            Går første år design, og har vært aktiv i økonomigruppen.

            \textbf{Sima Barzinje Velges ved akklamasjon.}

        \subsubsection{Innkjøpsansvarlig}
            % \textit{Vervet presenteres.}
            En kandidat til vervet.

            Mick
            Går fjerde året bachelor. Har vært aktiv på Cyb i et år. Har vært kaffemester i to semestere. Har drukket seg gjennom sortimentet allerede, og er lei. Ønsker fornyet sortiment.

            \textbf{Mick Williams velges ved akklamasjon.}

        \subsubsection{Teknisk sjef}
            % \textit{Vervet presenteres.}
            En kandidat til vervet

            Elias
            Er teknisk sjef nå.

            \textbf{Elias Hoel Birketvedt velges ved akklamasjon.}

        \subsubsection{Utlånsansvarlig}
            % \textit{Vervet presenteres.}

            En kandidat til vervet

            Silje
            Er utlånsansvarlig nå, og ønsker å fortsette.

            \textbf{Silje Helgesen velges ved akklamasjon.}


        \subsubsection{DJ-sjef}
            % \textit{Vervet presenteres.}
            En kandidat til vervet

            Peter
            Er nåværende DJ-sjef, og øsnker åfortsette i vervet.

            \textbf{Peter Tsunami velges ved akkklamasjon.}


        \subsubsection{Arrangementskoordinator}
            % \textit{Vervet presenteres.}

            En kandidat til vervet.

            Lina
            Er nåværende arrangementskoordinator, og synes det var gøy. Tror hun kommer til å ha tid.

            Miura
            Første året design. Er ikke tyve, men synes dette høres veldig gøy ut. Har jobbet her dette semesteret, og trives veldig godt. Ønsker å bidra mere.

            \textbf{Lina Hauck er valgt ved ordinært flertall.}

\end{multicols}



\section{Vedtektsforslag}

    Ingen innsendte forslag


\section{Utdeling av utmerkelser}

Utmerkelser utdeles.


\section{Eventuelt}

    Det er et definisjonsspørsmål om hva lovteksten "Hovedstyret i hende" skal bety når det kommer til vedtektsforslag. Skal det bety innsendt gjennom publiserte kanaler, eller skal det være nok at Hovedstyret er klare over forslaget.


\newpage

\textbf{\huge{Hovedstyret består nå av:}}

\mbox{}


\begin{center}
\setstretch{2.0}
    \LARGE Leder -- Peter Hjelle Petersen-Øverleir
    \\ Nestleder -- Philip Høegh Eriksen
    \\ Økonomiansvarlig -- Trym Shalchian-Tabrizi Auren
    \\ Kjellermogul -- Hanne Alexandersdatter Hugem
    \\ Arrangementssjef -- Eric Emanuel Reber
    \\ Rekrutteringsansvarlig -- Fredrik Arman Sarai
    \\ Internansvarlig -- Kristian Tereshchenko Stillingen
    \\ Promoteringssjef -- Ka Thas
    \\ X-sjef -- Kamil K\MakeLowercase{RZYŻANOWSKI}
\end{center}

\mbox{}
\\[24pt]

\begin{center}
\LARGE Protokollunderskrivere:\\[30pt]

\begin{multicols}{2}
    \Large
    Martin Kristoffer Jørstad Heggem\\[48pt]
    \rule{220pt}{1pt}

    Eric Emanuel Reber\\[48pt]
    \rule{220pt}{1pt}


\end{multicols}
\end{center}

% \addcontentsline{toc}{part}{Vedlegg}

% \newpage
% \phantomsection{}
% \addcontentsline{toc}{section}{Vedtekter for Cybernetisk Selskab} % chktex
% \includepdf[pages=-]{cyb_vedtekter.pdf}

\end{document}
